\section{多边形映射的马斯洛夫指标}
\label{sec:maslov-indices-polygonal-maps}

\stfhSubIndex{I125}{Masiluofuzhibiao}{马斯洛夫指标（Maslov index）}{Duobianxingyingshedezhibiao}{多边形映射的指标（of polygonal map）}
本节依照 (FOOO10a) 中的叙述，为\termPolygonalMap{}指派一个\termTopologicalIndex{}。

首先说明如何利用命题 3.2.14，把一个自然丛对关联给每个附着于\termTransversalLagrangianChain{}的\emph{多边形映射}。

考虑带有一组\termMarkedPoint{}
\[
  \{u_0,u_1,\ldots,u_k\}\subset\partial D^2
\]
的圆盘 $\Sigma=D^2$。把（边界）\termPuncturedDisc{}记为
\[
  \dot\Sigma=D^2\setminus\{u_i\}_{0\le i\le k}.
\]
对 $i=0,\ldots,k$，固定 $u_i$ 在 $\Sigma$ 中的一个小闭邻域
$U_i$，使各 $U_i$ 两两不交\footnote{译者注：原书写成“使各点
$u_i$ 互不相交”；这些标记点本来已经互异，而这里需要选取为两两不交的是
邻域 $U_i$。译文依此补正。}，并在每个邻域上配备一个\termConformalIsomorphism{}
\[
  U_i\setminus\{u_i\}\longrightarrow(-\infty,0]\times[0,1].
\]

令
\[
  \mathcal L=(L_0,L_1,\ldots,L_k)
\]
是拉格朗日子流形 $L_i$（$i=0,\ldots,k$）的一个有序
$(k+1)$-元组。

\begin{definition}
\label{def:transversal-lagrangian-chain}
若 $\mathcal L$ 两两横截，并且对任意互异的三元组 $i,j,\ell$ 都有
\[
  L_i\cap L_j\cap L_\ell=\varnothing,
\]
则称 $\mathcal L$ 为\termTransversalChain{}。称 $k+1$ 为该链的长度。
\end{definition}

现在考虑一条交点链
\[
  \vec p
  =\{p_{i(i+1)}\in L_i\cap L_{i+1}\mid i=0,\ldots,k\}
\]
以及一个光滑映射 $w:\dot\Sigma\to M$；它连续延拓到 $\Sigma$，并满足
\begin{equation}
  \lim_{z\to u_i}w(z)=p_{i(i+1)},
  \qquad
  w\bigl(\overline{u_i u_{i+1}}\bigr)\subset L_{i+1}.
\label{eq:polygonal-map-boundary-data}
\end{equation}
这里及下文的下标按模 $k+1$ 理解，即
$u_{k+1}=u_0$、$L_{k+1}=L_0$。\footnote{译者注：原书在
$i=k$ 时仍使用 $u_{i+1}$、$L_{i+1}$，却没有声明循环约定；译文将其明示。
此外，原书在\eqref{eq:polygonal-map-boundary-data}的第二个条件中写的是
$L_i$。但它令 $u_i$ 映到 $L_i\cap L_{i+1}$，且下一页起始终把 $u_i$
两侧的边界段分别配给 $L_i$ 与 $L_{i+1}$；译文依这套后续构造把此处补正为
$L_{i+1}$。}

\stfhIndex{I022}{00-pi2-l-p}{\stfhIndexPiTwoLMarkedP}
以
\[
  \pi_2(\mathcal L;\vec p)
\]
表示这类连续映射的同伦类集合，并以 $B$ 表示
$\pi_2(\mathcal L;\vec p)$ 的一个元素。接下来要相对于数据
$\mathcal L$、$\vec p$ 与 $B$，为这类映射 $w$ 指派一个拓扑指标。
考虑 $w^*TM$ 的一个\termSymplecticTrivialization{}：
\[
  w^*TM\cong
  \dot\Sigma\times(\mathbb R^{2n},\omega_0).
\]

在这个平凡化下，映射
\[
  \alpha_w:S^1=\partial\Sigma\longrightarrow\Lambda(n);
  \qquad
  t\longmapsto T_{w(t)}L_i
  \quad\text{若 }t\in\overline{u_{i-1}u_i}
\]
定义了一条\termPiecewiseSmoothPath{}；它在 $k+1$ 个点
$u_i\in\partial D^2$（$i=0,1,\ldots,k$）处不连续。\footnote{译者注：
原书把该映射的定义域写成 $S^1=\dot\Sigma$，但 $\dot\Sigma$ 是二维穿孔圆盘；
从分段条件 $t\in\overline{u_{i-1}u_i}$ 以及紧接着列出的不连续点可知，定义域应为
边界圆周 $\partial\Sigma$。译文依此补正。}
由横截性假设，$(\mathbb R^{2n},\omega_0)$ 中的两个拉格朗日子空间
$T_{w(u_i)}L_i$ 与 $T_{w(u_i)}L_{i+1}$\footnote{译者注：原书在此处的
第二个切空间以及下方路径端点中都把基点写成 $w(u_{i+1})$。然而横截对
$L_i,L_{i+1}$ 的共同交点是 $w(u_i)=p_{i(i+1)}$，而且随后选取的\personDarboux{}坐标卡也以该点为中心；译文两处均统一为 $w(u_i)$。}满足
\[
  T_{w(u_i)}L_i\cap T_{w(u_i)}L_{i+1}=\{0\}.
\]
从这些数据出发，我们将在适当构造的\termBorderedRiemannSurface{}
$\widehat{\dot\Sigma}$ 上关联出光滑\termSymplecticBundlePair{}
$(E_w,\lambda_w)$。

首先，对每个 $i=0,\ldots,k$，考虑辛向量空间
\[
  (S,\Omega)=\bigl(T_{w(u_i)}M,\omega_{w(u_i)}\bigr),
\]
并按照命题 3.2.14 选定一条光滑路径
\[
  \alpha_{i(i+1)}:[0,1]\longrightarrow\mathcal L(S,\Omega),
\]
使得
\[
  \alpha_{i(i+1)}(0)=T_{w(u_i)}L_i,
  \qquad
  \alpha_{i(i+1)}(1)=T_{w(u_i)}L_{i+1}.
\]

在每个 $p_{i(i+1)}\in L_i\cap L_{i+1}$ 处，都存在定义在
$V_i$ 上的\termDarbouxChart{}
\[
  \Phi_i:V_i\longrightarrow T_{p_{i(i+1)}}M,
\]
其中 $V_i$ 是 $p_{i(i+1)}$ 在 $M$ 中的一个邻域，并且\footnote{译者注：
原书在条件 (1) 以及条件 (3) 两个等式的右端把 $\Phi_i$ 误写成了 $\Phi$；
译文补回下标 $i$。}
\begin{enumerate}[label=(\arabic*)]
  \item $\Phi_i(p_{i(i+1)})=0\in T_{p_{i(i+1)}}M$；
  \item
  \[
    d\Phi_i(p_{i(i+1)}):
    T_{p_{i(i+1)}}M\longrightarrow T_{p_{i(i+1)}}M
  \]
  是恒等映射；
  \item
  \[
  \begin{aligned}
    \Phi_i(V_i\cap L_i)
      &=\Phi_i(V_i)\cap T_{p_{i(i+1)}}L_i,\\
    \Phi_i(V_i\cap L_{i+1})
      &=\Phi_i(V_i)\cap T_{p_{i(i+1)}}L_{i+1}.
  \end{aligned}
  \]
\end{enumerate}

\begin{exercise}
\label{ex:darboux-chart-transversal-pair}
证明这样的\personDarboux{}坐标卡存在。
\end{exercise}

于是，由\eqref{eq:polygonal-map-boundary-data}中的收敛性，可以把
$\dot\Sigma$ 完备化为紧致带边\personRiemann{}曲面\footnote{译者注：原书称非紧的
$\dot\Sigma$ 为“紧致带边\personRiemann{}曲面”，又说把它“分解”为随后才定义的
$\widehat{\dot\Sigma}$，并在端部等同中交替使用 $U_i$ 与
$U_i\setminus\{u_i\}$。译文依本段的构造改写为把 $\dot\Sigma$ 完备化，并在涉及
半无限带时明确删去穿孔。}
\[
  \widehat{\dot\Sigma}
  :=\dot\Sigma\cup\bigcup_i\widehat U_i,
\]
其中 $w(U_i)\subset V_i$。令
\[
  \widehat U_i\cong[-\infty,0]\times[0,1]
\]
为
\[
  (-\infty,0]\times[0,1]
  \cong U_i\setminus\{u_i\}
\]
的自然完备化（$i=0,\ldots,k$）。称这个并为 $\dot\Sigma$ 沿其各穿孔、
关于端部共形坐标卡
\[
  U_i\setminus\{u_i\}\cong(-\infty,0]\times[0,1]
\]
\stfhIndex{I164}{Shibaopo}{实爆破（real blow-up）}
所作的\emph{\termRealBlowUp{}}。

对 $p_{i(i+1)}$ 的邻域上给定的\personDarboux{}坐标卡 $\Phi_i$，映射
$\Phi_i\circ w$ 满足
\begin{equation}
  d\Phi_i\bigl(T_{w(\tau,0)}L_i\bigr)
    \equiv T_{p_{i(i+1)}}L_i,
  \qquad
  d\Phi_i\bigl(T_{w(\tau,1)}L_{i+1}\bigr)
    \equiv T_{p_{i(i+1)}}L_{i+1}.
\label{eq:darboux-end-trivialization}
\end{equation}
而且\footnote{译者注：原书在这里写成
$\Phi_i\circ w(\infty,t)\equiv p_{i(i+1)}$；但本段采用的是
$[-\infty,0]\times[0,1]$，且 $\Phi_i(p_{i(i+1)})=0$。译文相应补正端部符号与
坐标值。}
\[
  \Phi_i\circ w(-\infty,t)\equiv0.
\]
把 $w$ 到 $\widehat{\dot\Sigma}$ 的连续延拓记为
\[
  \widehat w:\widehat{\dot\Sigma}\longrightarrow M.
\]

\begin{lemma}
\label{lem:polygonal-pullback-extension}
$\dot\Sigma$ 上的拉回丛\footnote{译者注：原书在引理陈述中把主丛写成
$(\Phi\circ\widehat w)^*M$，把边界丛写成
$(\partial_j(\Phi\circ\widehat w))^*TL_j$，证明中又把 $w^*TM$ 写成
$w^*M$；这些写法混淆了拉回丛、\personDarboux{}坐标映射及其导数。原证明明确说平凡化由
$d\Phi_i$ 给出，译文据此统一为相应的切丛拉回，并以
$\widehat U_i$ 表示包含新增端边的完备化邻域。}$w^*TM$ 可以光滑延拓为
$\widehat w^*TM$。在 $\widehat U_i$ 上，这一延拓具有典范平凡化
\[
\begin{aligned}
  &\bigl(\widehat w^*TM;
    (\partial_i\widehat w)^*TL_i,
    (\partial_{i+1}\widehat w)^*TL_{i+1}\bigr)
    \big|_{\widehat U_i}\\
  &\qquad\cong
  (\widehat U_i;\partial_0\widehat U_i,\partial_1\widehat U_i)
  \times
  \bigl(T_{p_{i(i+1)}}M;
    T_{p_{i(i+1)}}L_i,T_{p_{i(i+1)}}L_{i+1}\bigr),
\end{aligned}
\]
其中 $\partial_i\widehat w$ 与 $\partial_{i+1}\widehat w$ 表示
$\widehat w$ 在 $\widehat U_i$ 两条边界边上的相应限制。
\end{lemma}

\begin{proof}
由\eqref{eq:darboux-end-trivialization}，$w^*TM$ 有一个平凡化，使得三元组
\[
  \bigl(w^*TM;(\partial_iw)^*TL_i,
    (\partial_{i+1}w)^*TL_{i+1}\bigr)
\]
限制到 $U_i\setminus\{u_i\}$ 后，在导数 $d\Phi_i$ 下映到平凡丛
\[
  (U_i\setminus\{u_i\};\partial_0U_i,\partial_1U_i)
  \times
  \bigl(T_{p_{i(i+1)}}M;
    T_{p_{i(i+1)}}L_i,T_{p_{i(i+1)}}L_{i+1}\bigr).
\]
引理立即由此得证。
\end{proof}

现在，我们已把拉回丛
\[
  w^*TM\longrightarrow\dot\Sigma
\]
光滑延拓为
\[
  \widehat w^*TM\longrightarrow\widehat{\dot\Sigma}.
\]
把相应的辛向量丛记为
\[
  E_w=\widehat w^*TM.
\]

接下来描述一个拉格朗日子丛\footnote{译者注：原书在这一段的三次边界记号中分别写成
$\partial\widehat{\dot\Sigma}$、$\partial\widehat\Sigma$ 以及
$\partial\dot{\widehat\Sigma}$；译文统一为本段已经构造的
$\partial\widehat{\dot\Sigma}$。}
\[
  \lambda\subset E_w|_{\partial\widehat{\dot\Sigma}}.
\]
在 $\dot\Sigma$ 的穿孔 $u_i$ 处，把路径
\[
  \alpha_{i(i+1)}:[0,1]\longrightarrow
  \Lambda\bigl(T_{w(u_i)}M\bigr)
\]
插入 $\{-\infty\}\times[0,1]$，放在两个\termGaussMap{}
\[
  \alpha_w|_{\overline{u_{i-1}u_i}}
  \quad\text{与}\quad
  \alpha_w|_{\overline{u_i u_{i+1}}}
\]
之间；这样就得到 $E_w|_{\partial\widehat{\dot\Sigma}}$ 的一个\termLagrangianSubbundle{}
\[
  \lambda_w\longrightarrow\partial\widehat{\dot\Sigma}.
\]

\stfhSubIndex{I126}{Masiluofuzhibiao}{马斯洛夫指标（Maslov index）}{DuobianxingMasiluofuzhibiao}{多边形马斯洛夫指标（polygonal Maslov index）}
\begin{definition}
\label{def:polygonal-maslov-index}
定义\termPolygonalMaslovIndex{}，记作 $\mu(w;\mathcal L,\vec p)$，为马斯洛夫指标
\[
  \mu(w;\mathcal L,\vec p)=\mu(E_w,\lambda_w).
\]
\end{definition}

当拉格朗日链 $\mathcal L$ 的长度为二，即 $k=1$ 时，这个指标以一种自然方式
与拉格朗日对子 $(L_0,L_1)$ 的\termMaslovViterboIndex{} (Vi88) 相关；我们将在
定义 13.6.2 中给出这种关系。

\begin{exercise}
\label{ex:maslov-viterbo-relation}
找出马斯洛夫指标与\personMaslov{}--\personViterbo{}指标之间的确切关系。
\end{exercise}

\begin{proposition}
\label{prop:polygonal-maslov-homotopy-invariance}
$\mu(w;\mathcal L,\vec p)$ 只依赖于同伦类
$[w]\in\pi_2(\mathcal L,\vec p)$ 以及路径 $\alpha_{i(i+1)}$ 的选取。
对 $B=[w]$，把相应的共同指标记为
$\mu(\mathcal L,\vec p;B)$。
\end{proposition}

\begin{proof}
设 $w_0,w_1$ 是定义在
\[
  \dot\Sigma=D^2\setminus\{u_i\}_{i=0,\ldots,k}
\]
上、附着于数据 $(\mathcal L,\vec p)$ 的两个相互同伦的多边形映射
\footnote{译者注：原书把这组穿孔的下标范围排成
$i=0,\mathord\cdot,k$；译文按本节此前的范围补为 $i=0,\ldots,k$。}
。它们之间的光滑同伦 $\{w_s\}_{0\le s\le1}$ 诱导一族辛向量丛
$w_s^*TM$。由于 $[0,1]$ 紧，可以找到 $u_i$ 的一个公共邻域 $U_i$，使
$w_s^*TM$ 被 $\Phi$ 平凡化，并且对所有 $s\in[0,1]$ 都有
\[
\begin{aligned}
  d\Phi:
  &\bigl(w_s^*TM|_{U_i};
    (\partial w_s)^*TL_i,(\partial w_s)^*TL_{i+1}\bigr)\\
  &\quad\longrightarrow
  (U_i;\partial_0U_i,\partial_1U_i)
  \times
  \bigl(T_{w(u_i)}M;T_{w(u_i)}L_i,T_{w(u_i)}L_{i+1}\bigr).
\end{aligned}
\]
我们用同一条路径 $\alpha_{i(i+1)}$ 填入线段\footnote{译者注：原书此处改写为
$\{\infty_i\}\times[0,1]\subset[0,\infty]\times[0,1]$，与本节此前固定的
$[-\infty,0]\times[0,1]$ 端部坐标相反；译文统一回负无穷端。}
\[
  \{-\infty_i\}\times[0,1]
  \subset[-\infty,0]\times[0,1]
  \cong\widehat U_i.
\]
再把上面的平凡化 $d\Phi$ 延拓到 $D^2$ 的其余部分，便得到丛对
$(E_{w_s},\lambda_{w_s})$ 的一个同伦。由丛对马斯洛夫指标的同伦不变性，得到
\[
  \mu(E_0,\lambda_0)=\mu(E_1,\lambda_1),
\]
证明完成。
\end{proof}

这个拓扑指标是计算构造\termFukayaCategory{}时出现的
$\overline\partial$-算子之线性化的\termFredholmIndex{}的基础。（参见
(FOOO10a)。）
